% ==================================================
% BACKGROUND
% ==================================================

\section{Background}

\subsection{Recommendation Systems Fundamentals}

Recommendation systems are computational methods that predict user preferences or ratings for items they have not yet interacted with. They are widely used in e-commerce, streaming platforms, social media, and increasingly in educational settings.

\subsubsection{Collaborative Filtering}

Collaborative filtering is based on the assumption that users with similar preferences in the past will have similar preferences in the future. It includes:

\begin{itemize}
    \item \textbf{User-based:} Find similar users and recommend items they liked
    \item \textbf{Item-based:} Find similar items and recommend based on user history
\end{itemize}

\subsubsection{Content-Based Filtering}

Content-based approaches recommend items similar to those the user has liked before, using item features and user profiles.

\subsection{Sequential Recommendation Models}

Sequential models consider the temporal ordering and dependencies in user interactions, which is crucial for domains like e-learning.

\subsubsection{Recurrent Neural Networks (RNNs)}

RNNs, including variants like LSTM and GRU, naturally capture sequential dependencies through their recurrent connections.

\subsubsection{Self-Attention Mechanisms}

Self-attention allows models to weigh the importance of different historical items when making recommendations, providing better interpretability than traditional RNNs.

\subsection{Graph-Based Sequential Recommendation}

Graph structures can represent relationships between users, items, and sequences. Graph neural networks (GNNs) combined with sequential models like gSASRec provide a powerful framework.

\subsubsection{The gSASRec Model}

The gSASRec (Graph Self-Attentive Sequential Recommendation) model combines:

\begin{enumerate}
    \item Graph neural network layers to capture item relationships
    \item Self-attention mechanisms to model temporal dependencies
    \item Sequence modeling to leverage historical interactions
\end{enumerate}

\subsection{The MARS E-Learning Dataset}

The MARS (Massive Open Online Courses) e-learning dataset contains:

\begin{itemize}
    \item Real student interaction logs from online courses
    \item Learning resource metadata and relationships
    \item Temporal information about when students access resources
    \item Multiple resource types (videos, quizzes, documents, etc.)
\end{itemize}

This dataset presents a realistic testbed for sequential recommendation research with domain-specific characteristics.

\newpage
